\documentclass[12pt]{article}

\usepackage{amsmath, amssymb, amsthm}
\usepackage{enumerate}
\usepackage{hyperref}
\usepackage{geometry}
\usepackage{newtxtext}
\usepackage{newtxmath}
\geometry{margin=1in}

\title{Graph Theory - Lecture 1}
\author{}
\date{}

% ----------------------------------------------------------------
\begin{document}
\maketitle

The AM-GM inequality is a very common tool in many an analyst's or combinatorialist's toolkit.
In this problem, we will see how it can be used to prove a very routine bound on edges in bipartite graphs.

% ================================================================
\section*{Questions}

\begin{enumerate}[(a)]

  \item Let $x$ and $y$ be nonnegative real numbers. Expand $(\sqrt{x} + \sqrt{y})^2$.
  \item Explain why $(\sqrt{x} + \sqrt{y})^2 \geq 0$.
  \item By shifting terms around, and doing it a little bit more than usual, show that $x + y \geq 2\sqrt{xy}$. Then, deduce the Arithmetic Mean-Geometric Mean inequality,
  which states that $\dfrac{x+y}{2} \geq \sqrt{xy}$. (Hint: I slipped a trap at the start. You may also need to redo exercise (a) and (b).)
  \item Show that $\left(\dfrac{x+y}{2}\right)^2 \geq xy$.
  \item A graph $X$ has a vertex set $V(X)$ and an edge set $E(X)$. A \textit{bipartite graph}
  is a graph whose vertex set can be partitioned into two sets $A$ and $B$,
  such that every edge has one endpoint in $A$ and one in $B$.
  We will show that if $X$ is a bipartite graph with $n$ vertices,
  then $|E(X)| \leq \dfrac{n^2}{4}$.
  \item Compute $|A| + |B|$ in terms of $n$.
  \item Write down the maximum number of edges between $A$ and $B$, in terms of $|A|$ and $|B|$.
  \item Compute $\dfrac{(|A|+|B|)^2}{4}$ for a bipartite graph with $n$ vertices.
  \item Using the previous parts, show that $|E(X)| \leq \dfrac{n^2}{4}$.

\end{enumerate}

% ================================================================
\newpage
\section*{Solutions}

Credit: Wikipedia and my teachers.

\begin{enumerate}[(a)]

  \item Expand by first in last out:
  \[
    (\sqrt{x}+\sqrt{y})^2 = x + 2\sqrt{xy} + y.
  \]

  \item Since $\sqrt{x}$ and $\sqrt{y}$ are real numbers, their sum $\sqrt{x}+\sqrt{y}$ is real,
  and any real number squared is nonnegative. Hence $(\sqrt{x}+\sqrt{y})^2 \geq 0$.

  \item From (b), $x + 2\sqrt{xy} + y \geq 0$, so
  \[
    x + y \geq -2\sqrt{xy}.
  \]
  Since both sides of $x+y \geq 2\sqrt{xy}$ follow from part (a):
  we have $(\sqrt{x}+\sqrt{y})^2 \geq 0$, which expands to $x+y+2\sqrt{xy}\geq 0$,
  but more directly, using the hint and checking quickly, $(\sqrt{x}-\sqrt{y})^2 \geq 0$ gives
  \[
    x - 2\sqrt{xy} + y \geq 0 \implies x+y \geq 2\sqrt{xy}.
  \]
  Dividing both sides by $2$ yields the AM-GM inequality:
  \[
    \frac{x+y}{2} \geq \sqrt{xy}.
  \]

  \item Squaring both sides of the AM-GM inequality (both sides are nonneg\-ative, so
  squaring preserves the inequality):
  \[
    \left(\frac{x+y}{2}\right)^2 \geq \left(\sqrt{xy}\right)^2 = xy.
  \]

  \item See the problem statement above.

  \item Since $\{A,B\}$ partitions $V(X)$, every vertex belongs to exactly one of $A$ or $B$, so
  \[
    |A| + |B| = n.
  \]

  \item Each vertex in $A$ can be connected to at most every vertex in $B$, and vice versa.
  Since edges only go between $A$ and $B$, the maximum number of edges is
  \[
    |A| \cdot |B|.
  \]

  \item Substituting $|A|+|B|=n$:
  \[
    \frac{(|A|+|B|)^2}{4} = \frac{n^2}{4}.
  \]

  \item From (g), $|E(X)| \leq |A|\cdot|B|$. From (d) applied with $x = |A|$ and $y = |B|$,
  \[
    |A|\cdot|B| \leq \left(\frac{|A|+|B|}{2}\right)^2 = \frac{n^2}{4}.
  \]
  Chaining these two inequalities gives $|E(X)| \leq \dfrac{n^2}{4}$. \qed

\end{enumerate}

\end{document}